\documentclass[11pt,a4paper]{article}
\usepackage[T1]{fontenc}
\usepackage{lmodern,microtype,amsmath,amssymb,graphicx,booktabs,tabularx,array,xcolor,fancyhdr}
\usepackage[margin=22mm,headheight=14pt]{geometry}
\usepackage[hidelinks,pdfauthor={},pdftitle={Thermoelectrically Cooled HgCdTe Reference Device: Engineering Assessment, Revision 3}]{hyperref}
\pdfinfoomitdate=1\pdftrailerid{}\pdfsuppressptexinfo=15
\definecolor{failure}{rgb}{.55,.12,.12}
\pagestyle{fancy}\fancyhf{}\fancyhead[L]{\small TEC-cooled HgCdTe reference device}\fancyhead[R]{\small Engineering assessment: Revision 3}\fancyfoot[C]{\thepage}
\setlength{\parindent}{0pt}\setlength{\parskip}{5pt}\setlength{\tabcolsep}{5pt}
\setcounter{tocdepth}{1}
\newcommand{\fig}[3]{\begin{figure}[!ht]\centering\includegraphics[width=#2\linewidth]{figures/#1.pdf}\caption{#3}\end{figure}}
\newcommand{\fail}{\textcolor{failure}{\textbf{FAIL}}}
\newcommand{\Die}{265.616093}
\newcommand{\Cold}{265.424850}
\newcommand{\Voltage}{0.551926633}
\newcommand{\GlobalResidual}{0.000144}
\newcommand{\ThermalResidual}{0.001838}
\newcommand{\CurrentResidual}{0.042206}
\newcommand{\WorstCut}{0.406120}
\newcommand{\BathRaw}{56.336914}
\newcommand{\BathWeak}{57.192663}
\newcommand{\ConvergenceRows}{A & 366,156 & 265.592524 & 0.551939871 & +41.348 & -47.828 \\
B & 458,061 & 265.613622 & 0.551945058 & +8.695 & -6.976 \\
C & 647,657 & 265.616093 & 0.551926633 & +6.477 & -4.609 \\}
\newcommand{\OperatingRows}{Low-parasitic vacuum & 230.975 & 1.15443 & 173.164 & 179.253 \\
Nominal Kr & 232.209 & 1.15193 & 172.790 & 182.196 \\
High-parasitic nitrogen & 235.031 & 1.14554 & 171.830 & 187.787 \\}
\newcommand{\ReducedDie}{232.209}
\newcommand{\ReducedVoltage}{1.15193}
\newcommand{\ReducedStageRows}{1 & 56.508 & 185.594 & 0.86057 & 76.029 \\
2 & 17.432 & 53.053 & 0.23748 & 21.241 \\
3 & 6.121 & 14.119 & 0.05332 & 4.921 \\}
\newcommand{\BudgetRows}{Terminal electrical input & +172.789624 \\
Imposed die heat & +2.000000 \\
Net header/cap/window heat into device & +9.021852 \\
Au warm-end conduction into die & +1.117526 \\
Cu pin-end thermal input (signed) & -2.732546 \\
Heat rejected at effective sink & +182.196456 \\}
\newcommand{\SpatialRows}{Transverse volume 1 & 784,696 & 265.616563 & 0.551929953 & -1.2885 \\
Transverse volume 2 & 1,099,293 & 265.616574 & 0.551930034 & -1.2907 \\}

\begin{document}
\thispagestyle{empty}
{\Large\bfseries Thermoelectrically Cooled HgCdTe Reference Device\par}
\vspace{4mm}{\large Engineering Assessment\par}
\vspace{4mm}{\large\bfseries Revision 3\par}
30 October 2026
\vspace{3mm}\hrule\vspace{4mm}
\textcolor{failure}{\textbf{Hypothetical reference article. Numerical release incomplete.}}

\textbf{Abstract.} A source-referenced nine-pin, two-ear housing and an independently selected three-stage cooler envelope define a synthetic HgCdTe detector article. This assessment combines two distinct numerical descriptions. A genuine coupled three-dimensional solid calculation resolves 88 thermoelectric legs, internal conductors, a header, explicit joint and attach layers, and the detector. At 0.1\,A, 300\,K underside bath and 2\,mW imposed die heat, its principal fine mesh gives approximately 265.6\,K die temperature and 0.552\,V. Current and conservative heat checks pass, while original raw boundary-gradient requirements remain unmet: the principal ceramic/attach mismatch is 6.48\% against 5\%, and the best independent normal-refinement bath error is 1.069\% against 1\%. Later transverse cases remain near 1.29\%. No passing combined mesh and systematically refined comparator have established numerical release.

A separate five-temperature model uses bounded nonlinear one-dimensional P/N transport, corrected wire routes, a configurable sink, imposed detector heating and preliminary gas/radiation exchange. Its enclosure is an explicitly filled stepped-core proxy inside candidate cap/window solids, not a fully resolved native cavity. At nominal 300\,K, 0.15\,A, 30\,kPa krypton and 2\,mW die heat, this model predicts approximately 232.2\,K and 1.152\,V. Of 63 prescribed scenario/current points, 52 are accepted numerically and 11 remain rejected with diagnostics. Material laws above 300\,K are assumed; no optical transmission, detector carrier transport or measured device-performance uncertainty is established.

\textbf{Engineering interpretation.} Source geometry, physical hypotheses, numerical verification and empirical accuracy are treated separately. Numerical digits identify conditional calculations, not manufacturing tolerances. Neither the VIGO housing plus candidate cooler reference nor the reduced enclosure model is a verified commercial assembly. The report supports bounded engineering exploration; it does not certify a device or a numerical release.
\clearpage
\tableofcontents
\clearpage
\section{Reference article, model hierarchy and interpretation}
The reference article is a three-stage thermoelectric cooler supporting a small HgCdTe die inside a nine-pin TO-66-style package. The intended questions concern detector temperature, electrical demand, heat rejection, parasitic paths and sensitivity to uncertain inputs. A proprietary assembly is not reconstructed, and agreement with a manufacturer's limiting numbers is not an acceptance requirement for this synthetic article.

\begin{tabularx}{\linewidth}{p{29mm}XX}\toprule
Description & Modeled physics and boundaries & Important exclusions\\\midrule
Three-dimensional solids & Coupled temperature and potential; 88 legs; 216 solid domains; 0.1\,A; 300\,K Kovar underside; 2\,mW die source & External Cu/Au links OFF; no cap, gas, radiation, optics or detector carriers\\
Reduced nonlinear package & Five unknown temperatures; nonlinear P/N profiles; Cu/Au conduction; Cu Joule heat; effective sink; detector heat & Ideal stage contacts and simplified bypass; no spatial spreading or die gradients\\
Candidate enclosure solids & Native cap/roof/window BREP and standalone 8,143-tetrahedron fixture & Geometry only; no conformal connection to the 88-leg mesh and no new thermal field\\
Preliminary exchange proxy & Closed cylinder and filled stepped core; gray radiosity; pressure-dependent gas chord network & Actual interleg voids, pin/wire occlusion, full native-surface view factors, gas diffusion/convection and spectral transmission absent\\\bottomrule
\end{tabularx}

\subsection{Input evidence and physical hypotheses}
\begin{tabularx}{\linewidth}{p{31mm}p{29mm}X}\toprule
Input family & Evidence category & Interpretation\\\midrule
Housing envelope/pin functions & Published drawing & Source-backed external dimensions; some contour details and seals inferred [1]\\
Cooler envelope & Independent datasheet & Candidate dimensions, not proof of assembly compatibility [2]\\
Stage split and routing & Inferred construction & 28/12/4 couples and internal paths are explicit hypotheses\\
P/N transport & Research analog & Measured specimen tables to 300\,K; assumed extension to 390\,K [4]\\
Joints, contacts, cap/window & Engineering assumptions & Explicit dimensions; unknown grades and contact impedances\\
Gas, pressure, emissivity, loads, sink & Operating scenarios & Selectable conditions, not observations or statistical priors\\\bottomrule
\end{tabularx}

\subsection{Meaning of verification}
A converged algebraic solve, current conservation, an accurate controlled benchmark and empirical device validation answer different questions. The first three can support a numerical implementation without identifying real materials or package construction. Raw-gradient qualification is stricter than a conservative reaction balance. Chosen low/nominal/high scenarios are engineering alternatives, not confidence intervals. The sections below keep these distinctions visible for every result.
\clearpage
\section{Geometry and source interpretation}
The governing package reference is the VIGO family drawing ZTM-TO66-Z030, Revision 9, dated 1 March 2024 [1]. Its two mounting ears and nine C pins of nominal diameter 0.76\,mm govern the corrected header mesh. Overall flange dimensions are 31.5 by 17.5\,mm, mounting-hole pitch 24.4\,mm and nominal hole diameter 3.7\,mm. The source does not uniquely establish every fillet, seal or internal dimension. Inferred contour radii and idealized glass seals remain explicit modeling assumptions.

The TEC Microsystems 3MC04-044-05 datasheet [2] independently provides a candidate three-stage cooler envelope: hot side 6.4\,mm, cold side 2.4\,mm and height 3.8\,mm. \textbf{The VIGO housing plus this cooler reference is not an identified OEM assembly.} Internal artwork, stage populations and material grades are unverified.

\fig{header_geometry}{.93}{Actual corrected tetrahedral mesh cut at header mid-height. Labels denote pin numbers in the modeled top-view coordinate system; the source bottom view is mirrored consistently. White mounting holes and the two-ear outline are mesh geometry, not an ellipse or an artist's reconstruction.}

\begin{tabularx}{\linewidth}{lX}\toprule
Element & Provenance and interpretation\\\midrule
Header and C pins & Drawing envelope; Kovar and glass thermal properties are proxies.\\
Pin functions & 1/9 TEC supply/return; 7/8 detector; 5/6 thermistor positions.\\
Internal TEC & 28/12/4 couples, totaling 44 pairs and 88 legs; assumed routing.\\
Detector and attach & 1\,mm$^2$ by 200\,$\mu$m die; 30\,$\mu$m attach; hypotheses.\\
Hot joint & Explicit 50\,$\mu$m metallic layer; unidentified alloy.\\\bottomrule
\end{tabularx}

\clearpage
\section{Actual solid topology and discretization}
The corrected mesh contains 216 solid domains. All 88 legs form one electrical series route through the copper pads and interstage conductors. There are two module terminals on the \emph{hot-side} copper pads. The external pin-to-pad copper and detector gold links are disabled in the present calculation: the terminal drive acts directly on the resolved module pad faces. The nine pins remain passive thermal solids with exposed ends adiabatic.

\fig{package_geometry}{.98}{Actual mesh exterior: two-ear header, nine feedthrough pins and the cooler/detector stack. No cap or window is modeled. Pin geometry is retained even though the physical external electrical links are OFF.}

The feedthrough seals and pin centers follow a separately audited coordinate registry. The header underside is the thermal bath; the extended pin shafts are not extra fixed-temperature boundaries. The view shows the solid domain used in the calculation, rather than a complete hermetic package.

\clearpage
\subsection{Mesh and conductor connectivity}
The solid meshes use Gmsh 4.13.1 geometry/meshing and Elmer 26.2.1 thermoelectric solution procedures. The principal field uses 159,991 nodes and 647,657 tetrahedra. Mesh checks identify material-pair contacts, exterior faces and terminal regions from actual tetrahedron adjacency, rather than inferring connectivity from a CAD rendering. No undeclared internal boundary is permitted. Every intended leg and interstage conductor belongs to the same two-terminal series component.

\fig{stack_mesh}{.96}{Exact $y=0.15$\,mm tetrahedron/plane intersection of the resolved stack. Void space remains white. Lines are actual cut-element edges; no scattered triangulation bridges gaps. The two hot-side module pad faces belong to the resolved copper routing. P legs are red, N legs blue, copper ochre, ceramics ivory, attach gray and die purple.}

\begin{tabularx}{\linewidth}{lX}\toprule
Resolved feature & Dimensions and interpretation\\\midrule
Ceramic plates & Widths 6.4, 4.55, 3.25 and 2.4\,mm; thickness 220\,$\mu$m each.\\
Leg sections & 0.4 by 0.4\,mm; stage heights 1.07, 0.82 and 0.64\,mm.\\
Stage allocation & 56, 24 and 8 legs from hot to cold; assumed internal design.\\
Current checks & 177 interior cuts across semiconductor and copper route.\\
Interfaces & Shared conformal nodes; finite bulk attach and joint layers.\\\bottomrule
\end{tabularx}

\clearpage
\section{Material laws, interfaces and uncertainties}
\begin{tabularx}{\linewidth}{lrrrX}\toprule
Solid & $k$ & $\sigma$ & $S$ & Qualification\\
 & W\,m$^{-1}$K$^{-1}$ & S\,m$^{-1}$ & $\mu$V\,K$^{-1}$ & \\\midrule
P leg & 1.5 & 140000 & $+160$ & Constant verification proxy\\
N leg & 1.5 & 165500 & $-170$ & Constant verification proxy\\
Copper & 400 & $5.814\times10^7$ & 0 & Constant proxy\\
Alumina & 30 & --- & 0 & Unidentified grade\\
Kovar/pins & 17.3 & --- & 0 & Assumed alloy response\\
Glass & 1 & --- & 0 & Seal proxy\\
Hot metallic joint & 35 & --- & 0 & Unidentified alloy\\
Die attach & 1 & --- & 0 & Unidentified material\\
HgCdTe & $k(T)$ & --- & 0 & Literature thermal model\\\bottomrule
\end{tabularx}

For Hg$_{0.71}$Cd$_{0.29}$Te, the literature model [3] is represented by piecewise linear nodal conductivity: $(T,k)=(90,2.41)$, $(148,1.61)$, $(298,0.973)$ and $(400,0.790)$ in kelvin and W\,m$^{-1}$K$^{-1}$. The calculation stays within this range. Electrical conductivity of passive domains is a computational extension of $10^{-12}$\,S\,m$^{-1}$, not a physical detector or pin current prediction. Density and heat capacity do not affect this steady calculation.

Interfaces have conformal shared nodes and ideal temperature/potential continuity. The explicit joint and attach have finite bulk thermal resistance; no additional calibrated contact impedance is added. Temperature-dependent, traceable P/N transport curves and measured contact resistances have \emph{not} been qualified for this hypothetical cooler. Earlier material studies cannot identify the actual grade of an unspecified assembly.

\subsection{Imposed conditions and exclusions}
The electrical drive is 0.1\,A with one hot-pad reference potential. Only the external Kovar underside is held at 300\,K. The die receives uniform volumetric heating totaling 2\,mW: 1.5\,mW imposed bias dissipation plus 0.5\,mW optical-like imposed heat. This does not simulate detector current, absorption, responsivity or transmission through a window.

All other exposed solid surfaces are adiabatic; no arbitrary 300\,K condition is imposed on the cold assembly or pins. Copper supply and gold detector hybrid links are OFF. The source drawing's cap and window, any carrier, thermistor internals, gas, radiation and carrier-transport PDEs are excluded. These exclusions apply to this solid-only field. The separate reduced package model later includes specified wire and enclosure exchanges without creating a three-dimensional cavity field.


\clearpage
\subsection{Temperature-dependent research analog for reduced transport}
The reduced model uses source-traced P/N research-specimen curves from Zhang et al.\ [4]. Complete measured-table coverage is 40--300\,K. The P and N samples are different laboratory specimens; they are not an identified commercial matched pair. The selected operating model restricts all leg temperatures to 180--390\,K. Between 180 and 300\,K, properties are interpolated on the source knots. Between 300 and 390\,K, the slope of each property is \emph{assumed} to be one quarter of its final measured-interval tangent. No measured validity above 300\,K is implied.

\fig{material_laws}{1.0}{P/N active material laws used only by the reduced analog models. Solid lines lie within the measured-table domain; dashed lines and the shaded region denote the assumed 300--390\,K continuation. The three-dimensional field shown elsewhere retains the constant verification properties.}

Resistivity $\rho$ is in $\Omega$\,m, conductivity is identically $\sigma=1/\rho$, thermopower $S$ is in V\,K$^{-1}$, and $k$ is in W\,m$^{-1}$K$^{-1}$. The Thomson coefficient is $\tau=T\,dS/dT$. At an interior source knot the right derivative is used; exactly 300\,K uses the measured left derivative, strictly above 300\,K the quarter tangent, and exactly 390\,K its left limit. A two-sided derivative is not asserted at the slope discontinuity.

All evaluated states must be finite, with positive $\rho$ and $k$ and the correct P/N sign of $S$. No solution is accepted by clipping it to a temperature limit. Interior Joule heating can activate the assumed continuation even when both boundary temperatures are below 300\,K. Continuation tests establish internal consistency only, not experimental validation. The research-specimen choice, digitization and assumed extension are consequential physical uncertainties.

\clearpage
\section{Coupled field equations and solution verification}
With electric potential $\phi$, $\mathbf E=-\nabla\phi$, conductivity $\sigma$, thermal conductivity $k$ and Seebeck coefficient $S$, the steady constitutive equations are
\begin{align}
\mathbf J&=\sigma(-\nabla\phi-S\nabla T), & \nabla\cdot\mathbf J&=0,\\
\mathbf q&=-k\nabla T+ST\mathbf J, & \nabla\cdot\mathbf q&=\mathbf J\cdot\mathbf E+q_v.
\end{align}
The coupled upstream Elmer thermoelectric solver includes Seebeck, Peltier, Joule and Fourier terms. It uses conformal linear tetrahedra and a nonlinear stopping threshold below $10^{-9}$. The reduced models introduced later are separate calculations; they do not replace or relabel this three-dimensional field.

\fig{temperature}{.88}{Actual solved temperature on the exact stack section. Colors represent source-cell means of the P1 field. The colder upper stage and warmer hot-side solids arise from the coupled solution. The field is conditional on the assumed materials and failed spatial qualification.}

\begin{tabularx}{\linewidth}{Xr}\toprule
Principal archived calculation & Numerical value\\\midrule
Nodes / tetrahedra & 159,991 / 647,657\\
Die volume mean & \Die\ K\\
Cold ceramic volume mean & \Cold\ K\\
Terminal voltage & \Voltage\ V\\
Known input, $I\Delta\phi+Q_{die}$ & 57.192663\,mW\\
Bath weak reaction / raw gradient & \BathWeak\ / \BathRaw\ mW\\\bottomrule
\end{tabularx}
Extra digits identify reproducible numerical outputs; they do not express physical accuracy or manufacturing tolerance.

\clearpage
\section{Electrical and heat-flow fields}
\fig{potential}{.81}{Actual electrical potential on the conductor section. Only P, N and copper are colored. The module terminals are on the hot side; external pin links are OFF.}
\fig{current}{.81}{Magnitude of the actual element-mean current density on the same section. Continuity is tested separately at every one of 177 interior cuts; plotting a smooth field alone is not a continuity test.}

\clearpage
\fig{heat}{.85}{Actual element-mean total solid heat flux, including Fourier and Peltier terms. The logarithmic scale reveals both the stack and passive conductive paths; it is not a recovered conservative boundary flux.}
\fig{cold_detail}{.98}{Actual resolved cold ceramic, 30\,$\mu$m attach and 200\,$\mu$m die section. The small die-temperature mesh sensitivity does not certify the normal heat-flux gradients. The maximum plotted range is a visualization choice, not an uncertainty band.}

\section{Independent raw-gradient diagnostics}
The heat-flux audit fits each tetrahedron's actual affine temperature gradient, uses its owner material conductivity and the oriented area of each physical triangle, and sums powers with area weights. At an internal interface, each solid supplies an independent outward trace. Define
\begin{equation}
e_\Gamma=\frac{Q^-_{out}+Q^+_{out}}{(|Q^-_{out}|+|Q^+_{out}|)/2},\qquad
e_b=\frac{Q_{b,raw}-(I\Delta\phi+Q_{die})}{I\Delta\phi+Q_{die}}.
\end{equation}
The unchanged limits are $|e_\Gamma|<5\%$ and $|e_b|<1\%$. Reactions and recovered fluxes are useful checks but cannot replace these raw criteria.

\clearpage
\section{Cold-interface and bath evidence}
\fig{cold_flux}{.93}{Actual independent raw Fourier flux on both sides of each cold contact, positive downward. The ceramic/attach mismatch is $+6.4765\%$ (FAIL); attach/MCT is $-4.6085\%$ (PASS). Different traces are numerical evidence, not an imposed contact jump.}

\fig{bath_flux}{.98}{Actual raw outward flux on the Kovar underside before and after a separate bath-volume refinement. The cold mesh is exactly fixed. Absolute integral errors improve from 1.4963\% to 1.2914\%, both FAIL. The reference is known external input, not the bath reaction.}

\clearpage
\section{Mesh sensitivity and independent analytical benchmark}
\begin{center}\small\begin{tabular}{lrrrrr}\toprule
Mesh & Tetrahedra & Die mean (K) & Voltage (V) & $e_{cer/att}$ (\%) & $e_{att/die}$ (\%)\\\midrule
\ConvergenceRows
\bottomrule\end{tabular}\end{center}
\fig{convergence}{.98}{Three actually completed coupled meshes. Successive die-temperature changes are 21.10 and 2.47\,mK; voltage is nonmonotone. Raw ceramic/attach error remains above 5\%. These are sensitivity data, not a rigorous whole-device grid-convergence index.}

The cold-region effective mesh sizes are 76.724, 25.484 and 17.608\,$\mu$m. Refinement is local and anisotropic; the bath and most TEC regions stay nearly fixed. A conditional temperature-only fit yields order 1.324 and a 3.91\,mK fine-to-limit difference, but this does \emph{not} qualify a device GCI. Nonmonotone voltage admits no such extrapolation. The separate 675,496-tetrahedron bath case shifts the die mean by $+0.421$\,mK and voltage by $+2.946\,\mu$V.

\clearpage
\subsection{Independent analytical reference}
The benchmark uses 220/30/200\,$\mu$m layers and conductivities 30/1/1.1\,W\,m$^{-1}$K$^{-1}$. Constant-flux and distributed-source cases separate normal-gradient truncation from integration errors. Coarser benchmark failures remain failures. Near-exact nodal values with finite mean and face-flux errors illustrate why an algebraically converged solution may still fail raw-gradient gates.

\fig{benchmark}{.98}{Independent three-dimensional, three-layer HeatSolver benchmark with an analytical temperature field, harmonic transverse variation and a distributed 2\,mW source. Source-face errors are 12.5, 6.25 and 3.125\%; mean-temperature errors are 1.855, 0.464 and 0.116\,mK. The finest benchmark passes its benchmark-specific preregistered checks; these do not replace the production bath criterion or qualify the coupled device.}

For a uniformly heated slab with constant conductivity, the axial solution includes a quadratic term, $T(z)=a+bz-q_vz^2/(2k)$. Its exact gradient varies linearly. A linear tetrahedron represents temperature by an affine field and therefore has a constant gradient. Even when nodal temperatures nearly match the analytical solution, the boundary trace samples a finite owner volume rather than the exact boundary derivative. Normal refinement addresses that truncation directly.

The analytical mean temperature is 300.1958788\,K. The finest distributed-source calculation gives 300.1957629\,K. Source-face flux error halves at each level, while mean-temperature error decreases by approximately four, consistent with first-order boundary gradients and second-order temperature integration in this controlled benchmark. The actual coupled device has material jumps, geometric edges and local refinement, so these rates are evidence for the mechanism, not transferable error estimates.

\clearpage
\section{Verification matrix and numerical release gates}
This matrix concerns the illustrated three-dimensional solid field. Independent refinement sequences are evaluated separately below; a reduced-model balance cannot certify a spatial flux gate.\par
\small
\begin{tabularx}{\linewidth}{p{37mm}p{34mm}Xl}\toprule
Test & Hard limit & Baseline evidence & Status\\\midrule
Source geometry/topology & 216 domains; 88 legs; two terminals & Corrected actual mesh audit & PASS\\
All 177 current cuts & Each $<0.5\%$ & Worst \WorstCut\% & PASS\\
Global heat residual & $<1\,\mu$W & \GlobalResidual\,$\mu$W & PASS\\
Free thermal residual & $<1\,\mu$W & \ThermalResidual\,$\mu$W & PASS\\
Free current residual & $<1\,\mu$A & \CurrentResidual\,$\mu$A & PASS\\
Nonlinear change & $<10^{-9}$ & $4.34\,10^{-10}$ & PASS\\
Homogeneous solid test & Passive 300\,K & Error below 4\,nK & PASS\\
Same-mesh replay & Agreement with native LU & 0.225\,$\mu$K; 3.70\,nV & PASS\\
Ceramic/attach raw flux & $<5\%$ absolute & 6.48\% & \fail\\
Attach/MCT raw flux & $<5\%$ absolute & 4.61\% & PASS\\
Independent raw bath & $<1\%$ absolute & Best normal case 1.069\% & \fail\\
Same production mesh and refined comparator & Both raw gates, all conservation checks & Not demonstrated & OPEN\\
Qualified whole-device GCI & Systematic asymptotic evidence & Local conditional fit only & OPEN\\
Empirical device accuracy & Independent measurements & No measured reference article & OPEN\\\bottomrule
\end{tabularx}
\normalsize

\textbf{Numerical release requires both original raw limits on one production mesh and at least one systematically refined comparator}, together with all 177 current cuts, global/stage heat checks, free thermal/electrical residuals, homogeneous baseline and reproducibility checks. Temperature and voltage sensitivity must be reported. Preferred 4.5\% cold and 0.9\% bath margins are targets, not relaxed or replacement gates. Passing two different meshes separately is insufficient.


\clearpage
\section{Independent normal and transverse refinement evidence}
The cold and bath refinement tracks retain the same 0.1\,A, 2\,mW, 300\,K underside boundary and disabled external links. Their equations and raw acceptance criteria are unchanged. Free owner-volume nodes are refined with region locks; changes to a temperature output do not substitute for flux checks.

\begin{center}\small\begin{tabular}{lrrrrr}\toprule
Study & Tetrahedra & Die mean (K) & Voltage (V) & Cold raw (\%) & Bath raw (\%)\\\midrule
Cold normal & 821,078 & 265.616118 & 0.551926634 & +4.795 & -1.496\\
Bath normal 1 & 744,713 & 265.616580 & 0.551930113 & +6.477 & -1.069\\
Bath normal 2 & 821,458 & 265.616593 & 0.551930210 & +6.477 & -1.093\\\bottomrule
\end{tabular}\end{center}
The cold-only track brings the ceramic/attach raw mismatch to 4.795\%, passing the original 5\% gate on that mesh; its bath still fails. The contact above it improves to 2.306\%. Die mean changes only 24.94\,$\mu$K while a boundary gradient changes materially. The bath-only normal-height sequence is approximately 296, 148 and 74\,$\mu$m; its 1.069\% then 1.093\% errors are nonmonotone, both FAIL. Separate tracks do not establish a same-mesh release.

\fig{latest_raw_flux}{.98}{Actual cold-normal ceramic trace and second bath-normal signed outward trace, from separate genuine solved fields. The bath map belongs to the 821,458-tetrahedron normal case. It is not the later transverse field and does not depict a combined passing mesh.}

\clearpage
\subsection{Stage conservation and later bath-volume behavior}
The second bath-normal field provides the following independent weak control-volume decomposition. The coldest band includes the 2\,mW die source, so its electrical-work difference alone is not the whole band balance.
\begin{center}\small\begin{tabular}{lrrrr}\toprule
Stage band & Hot out (mW) & Cold out (mW) & Electrical work (mW) & Residual ($\mu$W)\\\midrule
1 & 57.193021 & 18.014934 & 39.178087 & 0.000008\\
2 & 18.014934 & 5.754100 & 12.260833 & -0.000018\\
3 & 5.754100 & 0 & 3.754100 & -0.000006\\\bottomrule
\end{tabular}\end{center}
All three normal cases pass all 177 interior current cuts, global/stage heat and free-residual checks. The worst current-cut error is 0.406120\%. Maximum global heat error is 0.000093\,$\mu$W; maximum free thermal and current residuals are 0.002472\,$\mu$W and 0.042254\,$\mu$A. Those numerical checks remain separate from raw-gradient qualification.

Two additional transverse/volume bath refinements also solve the genuine coupled solid equations. Their cold/electrical regions remain fixed. Each passes the unchanged nonraw checks, but neither passes the original bath criterion:
\begin{center}\small\begin{tabular}{lrrrr}\toprule
Study & Tetrahedra & Die mean (K) & Voltage (V) & Bath raw (\%)\\\midrule
\SpatialRows
\bottomrule\end{tabular}\end{center}
The die-temperature difference between these two fields is approximately 11\,$\mu$K while bath errors plateau near 1.29\%. Signed local flux regions of opposite sign and internal gradient jumps matter; global cancellation can coexist with unresolved local traces. No recovered or reaction-based flux replaces the raw bath criterion. A combined cold-plus-bath passing production mesh and a systematically refined passing comparator are absent.

\subsection{Computational method and practical cost}
The fixed fine solid problem was solved with a native direct factorization and checked with a same-mesh numeric-refinement replay. The principal native calculation takes approximately 1,292\,s and 5.7\,GiB peak resident memory; a warm replay completes in approximately 558\,s and 5.8\,GiB. This is a warm-start comparison, not a general cold-start speed guarantee. The later coupled transverse fields complete in approximately 911 and 999\,s with three nonlinear iterations. These numerical costs motivate bounded transport/exchange reductions for engineering scenarios, without implying that the reduced model reproduces the full spatial boundary-value problem.

\clearpage
\section{Conservative nonlinear transport and five-temperature package}
The reduced package has unknown temperatures $T_h,T_1,T_2,T_c,T_d$: hot ceramic, two intermediate ceramics, cold ceramic and die. Each stage has a representative P leg and N leg, scaled by its assumed number of identical couples. No current-spreading field or per-leg manufacturing variability is solved.

Let $z$ run from hot to cold, $L$ be leg length and $A$ its area. Current densities are $j_N=I/A$ and $j_P=-I/A$. The steady one-dimensional laws are
\begin{align}
q&=STj-k\,dT/dz, & dq/dz&=j(\rho j+S\,dT/dz),\\
v_s&=j\rho(T_m)\Delta z+S(T_m)\Delta T, & w_s&=jv_s,\\
r_i&=q_{i+1/2}-q_{i-1/2}-\tfrac12(w_{i+1/2}+w_{i-1/2})=0.
\end{align}
Here $q$ is W\,m$^{-2}$, $v_s$ is V, $w_s$ and $r_i$ are W\,m$^{-2}$, $T_m$ is segment-midpoint temperature and $\Delta z=L/n$. Endpoint heat fluxes include the half-segment work. Both endpoint heats are independently computed:
\begin{equation}
\begin{aligned}Q_c&=-NA(q_{c,P}+q_{c,N})-K_b(T_h-T_c),\\
Q_h&=-NA(q_{h,P}+q_{h,N})-K_b(T_h-T_c),\end{aligned}\qquad V=N(v_N-v_P).
\end{equation}
$K_b$ is an explicit W\,K$^{-1}$ bypass hypothesis, zero in these scenarios. The identity $Q_h-Q_c=IV$ is audited rather than imposed by overwriting a flux. Variable thermopower already contributes through flux and electrical work; an extra Thomson heat source would double-count it.

An analytic tridiagonal Jacobian and banded Newton solve use a residual-decreasing bounded line search. Nominal legs have 32 intervals. Material domains are enforced at all evaluations. Previously evaluated stage problems are reused only for identical dimensions, boundary temperatures, current and material multipliers; inputs are not rounded. A constant-property package solution supplies an initial iterate only. Failed nonlinear/residual/domain checks remain rejected states with recorded candidate diagnostics.

\subsection{Package, wire and load balances}
At each intermediate node, preceding-stage $Q_c$ balances the next $Q_h$ plus net parasitic heat. The cold ceramic balances stage cooling against die-attach and local exchange; the die balances imposed heating, bond conduction and enclosure exchange against attach heat:
\begin{align}
Q_{c,0}&=Q_{h,1}+E_1, & Q_{c,1}&=Q_{h,2}+E_2,\\
Q_{c,2}&=E_c+G_a(T_d-T_c), & G_a(T_d-T_c)&=E_d+Q_{bias}+Q_{abs}.
\end{align}
$E$ is heat into the corresponding node in W; $G_a=k_aA_a/t_a$, nominally $1/30$\,W\,K$^{-1}$. The effective hot-side condition is
\begin{equation}
T_h-T_{amb}=R_{ha}\{Q_{h,0}+Q_{Cu,cond}+\tfrac12 I^2R_{Cu}+E_h\}.
\end{equation}
$R_{ha}=2$\,K\,W$^{-1}$ nominally represents the full hot-side path, not a resolved joint/header/spreader resistance. Corrected Cu routes join pins 1/9 to \emph{hot pads}; Au routes join pins 7/8 to the detector. Conductances are $kA/L$ and electrical resistances $\rho L/A$. Cu Joule heat is split between its warm pin and hot pad; terminal voltage includes that series loss. Bias and absorbed heat are prescribed, not inferred from an optical or carrier calculation.

\clearpage
\subsection{Independent transport accuracy and refinement}
Eight representative boundary-value cases keep all reference temperatures strictly inside the measured 40--300\,K interval. Two additional cases test the assumed high-temperature law across implementations only. The independent bounded nonlinear reference uses the same material tables and conservative midpoint equations but a different bounded least-squares algorithm. Same-discretization agreement checks implementation; 16/32/64 interval comparisons probe discretization, not physical material accuracy.

\fig{transport_accuracy}{1.0}{Absolute relative stage-voltage error against the independent 64-interval nonlinear reference. Stage numbering is hot to cold. The affine-temperature assumption omits Joule/Thomson-driven profile curvature; the direct nonlinear calculation substantially reduces that approximation error. These are numerical comparisons using research-specimen laws, not measurements of a device.}

At $T_h/T_c=260/240$\,K and 0.7\,A, nonlinear32 differs from the refined reference by 6.762\,$\mu$W in $Q_c$ (0.000823\%) and 0.1294\,mV (0.00570\%). The earlier affine voltage difference was about 147\,mV (6.48\%) against the 32-interval nonlinear reference. At 230/210\,K and 0.9\,A, nonlinear32 voltage error is 0.00823\%. The profiles stay bounded; interior temperature curvature, not a tuned coefficient, produces the improvement.

For the nominal package at 0.9\,A with the older effective-reservoir closure, nonlinear32-to64 changes die temperature by 0.000350\,K and terminal voltage by 0.1643\,mV. Knot crossings can cause small nonmonotone changes, so no whole-device asymptotic order or grid-convergence index is inferred. Median inner transport compute time over accepted scenario points is about 0.5\,s in the bounded four-CPU environment, excluding some initialization and outer-iteration overhead. Reusing identical stage evaluations reduces cost without introducing a surrogate interpolation error.

\clearpage
\section{Candidate enclosure and preliminary exchange geometry}
The candidate cap/roof/window are native solids using the documented cap outside diameter and explicit engineering hypotheses for wall, height and window dimensions. Cap height above the header is 8.60\,mm, wall 0.50\,mm, roof/window thickness 0.45\,mm and window diameter 5.90\,mm. The inner roof is at $z=9.65$\,mm, approximately 4.07\,mm above the die top. Steel-like cap and ZnSe-like window descriptions are candidate material concepts, not identified manufactured grades; no resolved window thermal/optical law is used.

\fig{candidate_enclosure}{1.0}{Left: central $y=0.15$\,mm sections of native candidate cap/window tetrahedra and the corrected solid stack, with geometric voids retained; the flange extends outside this view. The native enclosure is a separately meshed fixture, not a conformal addition to the coupled solid field. Right: the explicitly filled stepped-core geometry used only by the preliminary exchange operator. The native surfaces and proxy surfaces are different models. Neither panel is a newly solved enclosure temperature field.}

The native fixture has 8,143 tetrahedra and distinct shell, roof and window volumes. Positive tetrahedron volumes, unique exterior ownership, all nine physical face groups and shell/roof and roof/window interface incidence are audited. The package exterior is independently checked against the corrected solid mesh. Native exposed ceramic-top areas are 29.964, 16.1295, 8.9375 and 4.760\,mm$^2$; the die exposes approximately 0.955\,mm$^2$. These are not silently substituted for the different proxy areas.

The closed exchange geometry is an inner cylinder of radius 7.1\,mm with header floor and a filled stepped core based on the corrected widths and heights. Fourteen surface groups map to five unknown temperatures plus prescribed header, cap and window reservoirs. Vertical core groups use means of adjacent node temperatures. Actual interleg voids, exposed metallization, pin and wire occlusion are omitted. Thus the operator supplies a bounded, auditable preliminary load budget; full native-surface view factors remain unimplemented.

\clearpage
\subsection{Gray radiation and pressure-dependent gas network}
For gray opaque surface $i$, area $A_i$, emissivity $\epsilon_i$, radiosity $J_i$ and view factor $F_{ij}$,
\begin{align}
A_iF_{ij}&=A_jF_{ji}, & \sum_j F_{ij}&=1,\\
J_i&=\epsilon_i\sigma_{SB}T_i^4+(1-\epsilon_i)\sum_jF_{ij}J_j, &
Q_{rad,i}&=A_i\left(\sum_jF_{ij}J_j-J_i\right).
\end{align}
$\sigma_{SB}=5.670374419\,10^{-8}$\,W\,m$^{-2}$K$^{-4}$. Heat is positive into a surface. Deterministic scrambled low-discrepancy area/cosine-direction rays terminate on the first opaque proxy surface. Parallel upward faces, coplanar roof/window faces and convex self-views remain blocked. The concave wall can view itself. Raw exchange areas are symmetrized and symmetrically balanced for reciprocity and closure, retaining pairs unseen in both directions as zero. Both raw and corrected matrices are retained; conservation does not conceal the correction.

The gas path is a passive \emph{chord network}, not a solved diffusion or rarefied-gas field:
\begin{align}
G_{ij}&=\frac{k_g(T_{ij})A_iF_{ij}}{\ell_{ij}+2\delta_{ij}}, & Q_{g,i}&=\sum_jG_{ij}(T_j-T_i),\\
\delta_{ij}&=\frac{2-a}{a}\frac{2\gamma}{\gamma+1}\frac{\lambda}{Pr}, &
\lambda&=\frac{\mu}{p}\sqrt{\frac{\pi R_gT_{ij}}{2}},\\
\mu&=\frac{k_g}{c_p+1.25R_g}, & Pr&=\frac{\mu c_p}{k_g}.
\end{align}
$T_{ij}=(T_i+T_j)/2$ is the pair mean, $\gamma=c_p/c_v$ is the heat-capacity ratio, and $Pr$ is the dimensionless Prandtl number. $\ell_{ij}$ is a symmetric area-weighted ray chord in m, $G$ is W\,K$^{-1}$, $a$ is accommodation (assumed 0.8), $p$ is Pa, $R_g$ and $c_p$ are J\,kg$^{-1}$K$^{-1}$, and $\mu$ is Pa\,s. The viscosity relation is an Eucken engineering closure, not independent measured viscosity. At $p=0$, gas exchange is exactly zero; at small pressure it scales approximately with $p$. At large pressure it tends to the stated Fourier \emph{network} limit. Neither limiting check validates the network's physical geometry approximation.

The NIST gas conductivity tabulation [5] supplies N$_2$/Ar/Kr values at 200, 300 and 400\,K. At 300\,K, $k_g$ is 0.0259, 0.0177 and 0.0095\,W\,m$^{-1}$K$^{-1}$ respectively. Interpolation is bounded; 180--200\,K uses an explicitly flagged $T^{0.75}$ engineering continuation. No accurate accommodation, mixture/Xe transport, natural convection or production backfill optimum is established. The window is internally opaque gray; imposed absorbed heat is not a prediction of transmission.

\clearpage
\subsection{Ray refinement and conservative coupling}
\fig{ray_refinement}{1.0}{Bounded proxy-ray refinement at fixed package temperatures 300/270/240/220/221\,K and 300\,K reservoirs. Left: maximum absolute correction applied to a view factor during reciprocal closure. Right: resulting die gas and radiation heat input. These are deterministic refinement diagnostics, not statistical uncertainty bands or native-cavity error estimates.}

From 4,096 to 16,384 to 65,536 rays per surface, the largest absolute view-factor corrections are 0.07335, 0.05292 and 0.04043. The last value is still a 4.04-percentage-point correction. Die gas inputs are 0.136161, 0.139590 and 0.140812\,mW, and radiation inputs 0.289628, 0.293933 and 0.293511\,mW. Last-level changes are approximately 0.88\% gas and 0.14\% radiation. This does not qualify spatial view-factor accuracy; the filled-core model-form error is not estimated by ray refinement.

Every surface and effective-node exchange conserves heat to below $10^{-10}$\,W. Mapping surface heat back to package/reservoir nodes uses the same nonnegative temperature weights and preserves total power. Heat exchanged with the hot ceramic is included in the sink equation, not misapplied to the detector. Warm reservoirs are held at ambient; cap temperature gradients and their external cooling are not solved.

Independent two-plate gray-radiation, equal-temperature, zero-emissivity, vacuum, low-pressure scaling, high-pressure network, geometry-ownership and blocked-ray controls check implementation. Rejected material-domain and temperature-bound cases remain explicit. Exact reciprocal closure is a necessary consistency property, not proof of a physically accurate cavity operator.

\clearpage
\section{Operating scenarios, electrical demand and rejected states}
All reduced scenario families use 300\,K ambient and 2\,mW die heat. The low-parasitic family has vacuum, cap/window emissivities 0.08/0.05, 0.5\,K\,W$^{-1}$ sink and 0.75 times nominal Au conductance. The nominal family has 30\,kPa Kr, emissivities 0.2/0.1 and a 2\,K\,W$^{-1}$ sink. The high-parasitic family uses 101.325\,kPa N$_2$, emissivities 0.65/0.6, a 5\,K\,W$^{-1}$ sink, 1.25 times Au conductance and 1.5 times gas-network conductance. Several hypotheses change together; these families are not probability bounds.

\begin{center}\small\begin{tabular}{lrrrr}\toprule
At 0.15 A & Die (K) & Terminal (V) & Input (mW) & Sink (mW)\\\midrule
\OperatingRows
\bottomrule\end{tabular}\end{center}
These values use nonlinear 32-interval legs and the 65,536-ray filled-core exchange proxy. They are conditional engineering predictions. At the nominal point the hot ceramic is 300.364\,K, intermediate temperatures are 268.059 and 246.537\,K, and the cold ceramic is 232.104\,K. The die is warmer by approximately 0.105\,K because its attach carries both imposed and parasitic heat.

\fig{operating_sweep}{1.0}{Reduced nonlinear/proxy scenario sweeps from 0 to 1\,A in 0.05\,A increments. Missing intervals are numerical rejections, not interpolated operating data. Curves do not represent three-dimensional cavity solutions, measured performance or qualified continuous optima.}

Of 63 attempted points, 52 pass the reduced-model numerical gates and 11 remain rejected. Ten are lower-bound residual failures; the high-parasitic 1\,A attempt rejects a leg profile at its active temperature domain. A failed candidate does not prove that no physical operating state exists. The imposed 180\,K model floor also prevents interpreting accepted grid minima as unrestricted cooling optima. No source curve is silently extended below that floor.

\clearpage
\subsection{Material and closure comparisons at matched current}
At nominal 0.15\,A, the constant P/N proxy with the older effective-reservoir model gives die 243.280\,K and 0.84437\,V. The source analog with affine temperature gives 231.202\,K and 1.15539\,V. Direct nonlinear transport with the same older reservoir closure gives 231.192\,K and 1.15642\,V. Adding the current filled-core closure gives \ReducedDie\,K and \ReducedVoltage\,V. These alternatives change grade assumptions and/or package reduction; none establishes which hypothesis matches an unknown manufactured unit.

The approximately 12\,K constant-to-analog difference is far larger than the 0.010\,K affine-to-nonlinear change at this low-current point. At higher current the affine approximation can become materially inaccurate. The enclosure change at 0.15\,A raises the nominal die by approximately 1.018\,K; that shift includes changed surfaces, chords, reservoir paths and visibility assumptions, rather than a single independently measured physical mechanism.

The constant solid-only reduced case at 0.1\,A and 2\,mW gives 257.756\,K and 0.58561\,V when the hot ceramic is held at 300\,K. The genuine three-dimensional case holds only the Kovar underside at 300\,K and gives 265.593\,K and 0.55194\,V on its original mesh. Their difference is not a matched-boundary error estimate: spatial spreading, bypass and joint/header paths differ. It is retained as a warning about reduced-model structural omissions, not fitted away.

\subsection{Meaning of apparent numerical precision}
Several tables retain extra digits to distinguish solved states and audit balances. Without identified material grades, contact impedances, die loading, accommodation and reservoir conditions, such digits do not provide a physical accuracy interval. A microkelvin-scale change in an output can coexist with a failed raw spatial flux criterion, as the solid refinements demonstrate. A small reduced-model residual similarly cannot certify a real detector temperature.

\clearpage
\section{Heat budgets and dominant paths}
At nominal 0.15\,A, total die heat is 3.507\,mW: the 2\,mW imposed source plus approximately 1.507\,mW gas/radiation/Au input. The cold ceramic also absorbs approximately 2.614\,mW, so the coldest stage removes approximately 6.121\,mW. Interstage loads accumulate toward the hot stage.

\fig{reduced_heat_budget}{1.0}{Net positive heat loads at the four cold/intermediate reduced nodes for nominal 0.15\,A Kr. Attach transfer and interstage thermoelectric flows are internal and are not stacked again as external loads. Vertical proxy surfaces distribute heat using their stated node weights.}

\begin{center}\small\begin{tabular}{lrrrr}\toprule
Stage & $Q_c$ (mW) & $Q_h$ (mW) & $V$ (V) & Joule heat (mW)\\\midrule
\ReducedStageRows
\bottomrule\end{tabular}\end{center}
\begin{center}\small\begin{tabular}{lr}\toprule
Whole reduced device, signed inputs & Power (mW)\\\midrule
\BudgetRows
\bottomrule\end{tabular}\end{center}
The first five signed entries sum to sink rejection. Warm-reservoir power is explicitly accounted for; it is not part of terminal electrical input. The hot-side Cu connection carries heat outward because the pad is warmer than its ambient pin. The nominal global discrepancy is below $10^{-11}$\,W; the maximum over accepted proxy grid points is approximately $1.41\,10^{-11}$\,W. Those figures establish bookkeeping consistency, not a physical heat-leak uncertainty.

\clearpage
\section{Sensitivity, model-form uncertainty and physical interpretation}
\fig{sensitivity}{1.0}{One-factor engineering endpoint spans about nominal 0.15\,A, with nonlinear transport and the filled-core closure. Trial intervals are shown explicitly. They are chosen hypothesis ranges, not probability distributions or confidence intervals. Different trial widths would change the ranking.}

P/N thermal conductivity scaled from 0.8 to 1.2 times the selected law changes the die by 18.900\,K; Seebeck scaled from 0.9 to 1.1 changes it by 11.142\,K. Resistivity scaled 0.8--1.2 gives 2.371\,K, and bias heat from 0.5 to 5\,mW gives 2.226\,K. In this particular closure, vacuum to 101.325\,kPa Kr spans 0.721\,K, sink resistance 0.5--5\,K\,W$^{-1}$ spans 0.446\,K, and cap emissivity 0.08--0.65 spans 0.413\,K.

These results support prioritizing grade/contact and loading hypotheses when interpreting absolute cooling. They do not prove that real gas or radiation is negligible: the proxy omits interleg spaces, pin/wire occlusion, nonisothermal cap behavior and physical gas diffusion. Its geometry/model-form error is distinct from the reported ray-refinement changes. Ideal contacts and zero additional stage bypass also leave unresolved paths that can change both cooling and voltage.

Unknown construction and operating inputs should therefore be treated as alternative engineering hypotheses. Independent thermal/electrical measurements would constrain them, but absent proprietary data need not prevent bounded synthetic studies. Empirical comparison, when available, requires matched boundary temperature, current, imposed load and ambient conditions. No inferred confidence interval or commercial performance claim follows from the present sensitivity spans.

\clearpage
\section{Overall verification and engineering conclusions}
\small\begin{tabularx}{\linewidth}{p{35mm}p{30mm}X}\toprule
Assessment & Result & Interpretation\\\midrule
Three-dimensional geometry/current & PASS numerical checks & Corrected mesh and 177 interior current cuts\\
Three-dimensional weak heat/free residuals & PASS numerical checks & Algebraic and control-volume consistency\\
Original cold/bath raw criteria & \fail\ overall & Initial 6.48\% cold; best normal bath 1.069\%; later transverse about 1.29\%\\
Same-mesh release/comparator & Not demonstrated & Both raw limits and all other checks are still required\\
Bounded nonlinear leg transport & PASS tested comparisons & Same-grid/analytical/refined numerical checks; not grade validation\\
Reduced package conservation & PASS numerical checks & Accepted states meet unchanged energy/domain gates; rejected candidates retained\\
Proxy visibility/reciprocity/closure & PASS consistency checks & Closed filled-core rays and balanced exchange; 0.04043 maximum last-level correction\\
Native full-cavity exchange & Not implemented & No actual all-surface view factors, gas diffusion or new coupled cavity field\\
Material and device accuracy & Not established & Research analogs, assumed extensions and construction remain uncertain\\\bottomrule
\end{tabularx}\normalsize

\subsection{Engineering conclusions}
The source-referenced solid model demonstrates a two-terminal 88-leg coupled calculation with meaningful electrical/heat verification, while exposing unresolved raw boundary gradients. Stable temperature and voltage do not remove the original release failures. The direct nonlinear transport option markedly improves high-current consistency with a refined one-dimensional reference at practical scenario cost. Its physical grade and high-temperature continuation remain assumptions.

The preliminary enclosure operator makes reservoir heat exchange explicit and conservative, and records its geometric corrections. Its filled-core closure is useful for bounded comparisons, but it should not be interpreted as an accurate hermetic-cavity field. Native interleg/pin surfaces, contact and bypass paths, and cap reservoir behavior are the important unresolved physical descriptions. A large geometry refinement alone cannot establish their correctness.

For the specified nominal reduced scenario, approximately 232.2\,K die temperature, 0.173\,W electrical input and 0.182\,W sink rejection describe one conditional synthetic operating point. Material-law and imposed-load sensitivity dominate the chosen scenario spans. Neither an optimum current nor a defensible device-performance confidence interval is established. The assessment supports transparent engineering exploration while maintaining a clear boundary between computation, hypothesis and measurement.

\subsection{Reproducible scientific method}
The three-dimensional figures use exact intersections of stored tetrahedra and solved fields, with independent node/cell correspondence; voids are retained. The reduced figures are calculated from completed scenario data and independent boundary-value comparisons, with failures visible. Native candidate geometry is audited separately from the proxy. Sources, numerical inputs, archived field identities and deterministic figure/build records are retained outside the reader-facing pages. No field is invented to draw a figure and no additional production solve is implied by a figure rebuild.

\clearpage
\phantomsection
\addcontentsline{toc}{section}{References}
\begin{thebibliography}{9}
\bibitem{vigo} VIGO Photonics, \emph{3TE-TO66(9p)-wW PVI/PCI Detector Technical Drawing}, ZTM-TO66-Z030, Revision 9, 1 March 2024. Public family drawing; no identified composite assembly.
\bibitem{tec} TEC Microsystems, \emph{3MC04-044-05 Thermoelectric Cooler Datasheet}, 2019. Independent candidate cooler envelope.
\bibitem{mct} E. M. Heckman, L. P. Gonzalez and S. Guha, \emph{Measurement of optical and thermal properties of Hg$_{1-x}$Cd$_x$Te}, Applied Optics \textbf{47}, 578--582 (2008), \href{https://doi.org/10.1364/AO.47.000578}{doi:10.1364/AO.47.000578}. Composition-dependent thermal model.
\bibitem{zhang} J. Zhang et al., \emph{Accurate Measurements and Error Analysis of Bi$_2$Te$_3$-Based Low-Temperature Thermoelectrics}, \emph{Chinese Physics Letters} \textbf{42}, 056304 (2025), \href{https://doi.org/10.1088/0256-307X/42/5/056304}{doi:10.1088/0256-307X/42/5/056304}. Digitized P/N specimen transport, 40--300\,K; no commercial grade identification.
\bibitem{gas} M. L. Huber, \emph{Thermal Conductivity of Gases}, NIST/CRC (2019), \href{https://tsapps.nist.gov/publication/get_pdf.cfm?pub_id=926438}{NIST thermophysical-property tabulation}. Dilute-gas conductivity data; package transport is separately approximated.
\end{thebibliography}
\end{document}
